\documentclass[10pt,aspectratio=169]{beamer}
% Preámbulo modular para presentaciones Beamer (Modelación Estadística)
% Diseño y paleta institucional del Tecnológico de Monterrey

\documentclass[aspectratio=169,xcolor={dvipsnames,table}]{beamer}

\usepackage[utf8]{inputenc}
\usepackage[spanish,mexico]{babel}
\usepackage{amsmath,amssymb,amsfonts,mathtools}
\usepackage{graphicx}
\usepackage{listings}
\usepackage{booktabs}
\usepackage{tikz}
\usetikzlibrary{matrix,arrows,decorations.pathmorphing,shapes.geometric,positioning}

% Paleta Institucional Tecnológico de Monterrey
\definecolor{TecAzul}{HTML}{0039A6}       % Azul TEC: color primario institucional
\definecolor{TecAzulOscuro}{HTML}{1A2E51} % Azul marino oscuro
\definecolor{TecRojo}{HTML}{EC2661}       % Rojo de acento (paleta miTec)
\definecolor{TecGris}{HTML}{646464}       % Gris neutro para comentarios y subtítulos
\definecolor{TecFondoBloque}{HTML}{F4F6F9}% Fondo suave para bloques

% Configuración de Tema Beamer
\usetheme{Boadilla}
\usecolortheme[named=TecAzulOscuro]{structure}
\setbeamercolor{alerted text}{fg=TecRojo}
\setbeamercolor{palette primary}{bg=TecAzulOscuro,fg=white}
\setbeamercolor{palette secondary}{bg=TecAzul,fg=white}
\setbeamercolor{palette tertiary}{bg=TecAzulOscuro!80!black,fg=white}
\setbeamercolor{title}{fg=TecAzulOscuro,bg=TecFondoBloque}
\setbeamercolor{frametitle}{fg=TecAzulOscuro,bg=TecFondoBloque}

% Bloques personalizados elegantes
\setbeamercolor{block title}{bg=TecAzulOscuro,fg=white}
\setbeamercolor{block body}{bg=TecFondoBloque,fg=black}
\setbeamercolor{block title alerted}{bg=TecRojo,fg=white}
\setbeamercolor{block body alerted}{bg=TecRojo!8,fg=black}
\setbeamercolor{block title example}{bg=TecAzul,fg=white}
\setbeamercolor{block body example}{bg=TecAzul!8,fg=black}

% Redondear bordes y añadir sombra sutil
\setbeamertemplate{blocks}[rounded][shadow=true]
\setbeamertemplate{navigation symbols}{}

% Pie de página limpio con número de diapositiva
\setbeamertemplate{footline}{
  \leavevmode%
  \hbox{%
  \begin{beamercolorbox}[wd=.7\paperwidth,ht=2.25ex,dp=1ex,left]{author in head/foot}%
    \hspace*{2ex}\insertshortauthor~~(\insertshortinstitute) \hfill \insertshorttitle
  \end{beamercolorbox}%
  \begin{beamercolorbox}[wd=.3\paperwidth,ht=2.25ex,dp=1ex,right]{date in head/foot}%
    \insertshortdate{}\hspace*{2em}
    \insertframenumber{} / \inserttotalframenumber\hspace*{2ex}
  \end{beamercolorbox}}%
  \vskip0pt%
}

% Configuración de Listings de Python para visualización pedagógica de código
\lstset{
    language=Python,
    basicstyle=\ttfamily\footnotesize,
    keywordstyle=\color{TecAzulOscuro}\bfseries,
    stringstyle=\color{TecRojo},
    commentstyle=\color{TecGris}\itshape,
    showstringspaces=false,
    breaklines=true,
    frame=lines,
    rulecolor=\color{TecAzulOscuro!30},
    backgroundcolor=\color{TecFondoBloque},
    aboveskip=0.5em,
    belowskip=0.5em,
    literate={á}{{\'a}}1 {é}{{\'e}}1 {í}{{\'i}}1 {ó}{{\'o}}1 {ú}{{\'u}}1
             {Á}{{\'A}}1 {É}{{\'E}}1 {Í}{{\'I}}1 {Ó}{{\'O}}1 {Ú}{{\'U}}1
             {ñ}{{\~n}}1 {Ñ}{{\~N}}1 {¿}{{?`}}1 {¡}{{!`}}1
}

% Metadatos institucionales por defecto
\title[Modelación Estadística]{Modelación Estadística para Ciencia de Datos e Ingeniería}
\author[J. Castillo Colmenares]{Juliho David Castillo Colmenares}
\institute[Tec de Monterrey]{Tecnológico de Monterrey \\ Escuela de Ingeniería y Ciencias}
\date{\today}
% Comandos y macros estadísticas compartidas para las presentaciones Beamer (Metropolis)

% Conjuntos numéricos básicos
\newcommand{\R}{\mathbb{R}}
\newcommand{\N}{\mathbb{N}}
\newcommand{\Q}{\mathbb{Q}}
\newcommand{\Z}{\mathbb{Z}}
\newcommand{\set}[1]{\left\{ #1 \right\}}
\newcommand{\abs}[1]{\left|#1\right|}
\newcommand{\norm}[1]{\left\| #1 \right\|}

% Comandos estadísticos y probabilísticos del curso
\newcommand{\Var}{\operatorname{Var}}
\newcommand{\cov}{\operatorname{Cov}}
\newcommand{\corr}{\rho}
\newcommand{\comb}[2]{\begin{pmatrix} #1 \\ #2 \end{pmatrix}}
\newcommand{\s}{\sigma}
\newcommand{\particion}{\mathfrak{P}}
\newcommand{\card}[1]{\#\left( #1 \right)}
\newcommand{\seq}[3]{\set{#1}_{#2}^{#3}}
\newcommand{\E}{\mathbb{E}}
\newcommand{\Prob}{\mathbb{P}}

% Entornos pedagógicos y cajas en español académico natural
\newenvironment{discusion}{%
  \begin{alertblock}{Pregunta de Discusión}%
}{%
  \end{alertblock}%
}

\newenvironment{reflexion}{%
  \begin{alertblock}{Reflexión para el Aula}%
}{%
  \end{alertblock}%
}

\newenvironment{paradoja}{%
  \begin{alertblock}{Paradoja de Exploración}%
}{%
  \end{alertblock}%
}

\newenvironment{motivacion}{%
  \begin{exampleblock}{Motivación: Ciencia de Datos en la Práctica}%
}{%
  \end{exampleblock}%
}

\newenvironment{retoclase}{%
  \begin{block}{Reto Rápido en Equipo}%
}{%
  \end{block}%
}

\title[Diseños factoriales]{Sección 08.07: Introducción a diseños factoriales}\subtitle{Efectos principales e interacción}\author[J. Castillo Colmenares]{Juliho Castillo Colmenares}\institute{Tecnológico de Monterrey}\date{}
\begin{document}
\begin{frame}[plain]\titlepage\end{frame}
\begin{frame}{Hoja de ruta}\small\begin{enumerate}\item Diseño completo.\item Efectos principales e interacción.\item Modelo, ANOVA e interpretación.\end{enumerate}\end{frame}
\begin{frame}{Fuente y alcance}\small La sección extiende los diseños de un factor a dos o más factores de tratamiento observados simultáneamente.\begin{block}{Pregunta guía}¿Cómo detectar que el efecto de un factor depende de otro?\end{block}\end{frame}
\begin{frame}{Diseño factorial completo}\small Con factores $A$ y $B$, un diseño completo observa todas las $a\times b$ combinaciones, con $n$ réplicas por celda y $N=abn$.\end{frame}
\begin{frame}{Ventaja frente a OFAT}\small Variar un factor a la vez no permite estimar interacciones. El diseño factorial usa las mismas corridas para estudiar efectos principales y conjuntos.\end{frame}
\begin{frame}{Efecto principal}\small El efecto principal de $A$ es el cambio promedio de la respuesta al cambiar su nivel, promediando sobre los niveles de $B$. Análogamente se define el efecto de $B$.\end{frame}
\begin{frame}{Interacción}\small Hay interacción cuando el efecto de $A$ depende del nivel de $B$. En una gráfica, las curvas de respuesta no son paralelas.\end{frame}
\begin{frame}{Modelo factorial}\small\[Y_{ijk}=\mu+\alpha_i+\beta_j+(\alpha\beta)_{ij}+\epsilon_{ijk}.\]\end{frame}
\begin{frame}{Restricciones}\small Las restricciones de suma cero identifican los efectos: $\sum_i\alpha_i=\sum_j\beta_j=0$ y las sumas de interacción por filas y columnas también son cero.\end{frame}
\begin{frame}{Partición de ANOVA}\small\[\mathrm{SCT}=\mathrm{SC}_A+\mathrm{SC}_B+\mathrm{SC}_{AB}+\mathrm{SCE}.\]Cada fuente tiene su cuadrado medio y estadístico $F$.\end{frame}
\begin{frame}{Orden de interpretación}\small Interpretar de adentro hacia afuera: primero la interacción $F_{AB}$; solo si no es significativa se interpretan directamente los efectos principales.\end{frame}
\begin{frame}{Procedimiento I}\small\begin{enumerate}\item Definir factores, niveles y respuesta.\item Observar todas las combinaciones.\item Replicar y aleatorizar celdas.\end{enumerate}\end{frame}
\begin{frame}{Procedimiento II}\small\begin{enumerate}\setcounter{enumi}{3}\item Ajustar el modelo factorial.\item Evaluar interacción y efectos.\item Reportar medias por combinación.\end{enumerate}\end{frame}
\begin{frame}{Aplicación: algoritmos y lotes}\small Factor $A$: algoritmo ($A_1,A_2$). Factor $B$: tamaño de lote ($B_1=32,B_2=256$). Respuesta: tiempo de entrenamiento.\end{frame}
\begin{frame}{Aplicación: medias celulares}\small Las medias observadas son $42,40$ para $A_1$ y $38,20$ para $A_2$, al cambiar de lote pequeño a grande.\end{frame}
\begin{frame}{Aplicación: efectos simples}\small Bajo $A_1$, el cambio es $-2$ minutos; bajo $A_2$, es $-18$ minutos. El efecto del lote depende del algoritmo.\end{frame}
\begin{frame}{Aplicación: interacción}\small La diferencia entre efectos simples indica interacción sustancial. No es correcto afirmar que el lote grande reduce el tiempo en todos los algoritmos.\end{frame}
\begin{frame}{Aplicación: contraste}\small Para confirmar formalmente, evaluar $F_{AB}$ con $(a-1)(b-1)$ grados de libertad del numerador y $ab(n-1)$ del error.\end{frame}
\begin{frame}{Interpretación}\small Con interacción significativa, reportar combinaciones específicas. Los efectos principales aislados pueden ocultar direcciones opuestas.\end{frame}
\begin{frame}{Comunicación}\small Incluir tabla de medias por celda, tamaños, interacción, efectos principales y una gráfica de perfiles cuando sea posible.\end{frame}
\begin{frame}{Síntesis}\small\begin{itemize}\item El factorial cruza todos los niveles.\item La interacción expresa dependencia entre factores.\item El orden de interpretación empieza por $AB$.\end{itemize}\end{frame}
\begin{frame}{Conexión con el capítulo}\small Los diseños factoriales amplían el control experimental y preparan modelos con más factores e interacciones.\end{frame}
\end{document}
